\uuid{5XGQ}
\exo7id{4404}
\titre{exo7 4404}
\auteur{quercia}
\organisation{exo7}
\datecreate{2010-03-12}
\isIndication{false}
\isCorrection{true}
\chapitre{Série numérique}
\sousChapitre{Fonction exponentielle complexe}
\cpge{mpsi-mp2i.nombres-complexes.d, pcsi.nombres-complexes.d, ptsi.nombres-complexes.d, mpsi-mp2i.nombres-complexes.g, pcsi.nombres-complexes.g, ptsi.nombres-complexes.g, mpsi-mp2i.nombres-reels-suites.d, pcsi.nombres-reels-suites.c, ptsi.nombres-reels-suites.c, bcpst1.nombres-complexes.b, bcpst1.suites-reelles}
\module{Analyse}
\niveau{L1}
\difficulte{}

\contenu{
\question{
Soit $z\in\C$. Montrer que $\left(1+\frac zn\right)^n \to e^z$ lorsque $n\to\infty$.
}
\reponse{
Mettre $1+\frac zn$ sous forme trigonométrique.
}
}
